% !PRINTED 37
\weberpage{37}
One obtains further [without lengthy calculation, by comparison of the coefficients of the highest power of $\xi$ in (1)]
\[
\begin{aligned}
T &= \frac{a}{b'}(a\xi+b\eta)(a'\xi+b'\xi)(\alpha^2\xi-\beta^2\eta)(\beta^2\xi-\varkappa^2\alpha^2\eta), % [sic] second factor prints $\xi$ in both terms, no $\eta$
\\
H &= \frac{1}{3}\left(\frac{\partial U}{\partial \xi}\frac{\partial V}{\partial \eta}-\frac{\partial V}{\partial \xi}\frac{\partial U}{\partial \eta}\right) = \frac{a'}{a}T,
\end{aligned}
\]
from which one easily finds by (8):
\weq{13}{M = \frac{a}{a'} = \frac{\alpha}{\alpha+2\beta} = \frac{v}{v+2u^3}.}

If one substitutes the expression resulting from this,
\[
v = \frac{2u^3M}{1-M}
\]
into equation (12), then for $M$ there results an equation of the fourth degree, the \gesperrt{multiplier equation}, which can replace the modular equation. But one obtains this equation more simply in the following way. By (13),
\[
\frac{1}{M}-1 = 2\frac{\beta}{\alpha}, \quad \frac{1}{M}+3 = \frac{2(\beta+2\alpha)}{\alpha},
\]
hence by (9):
\[
\left(\frac{1}{M}-1\right)^3\left(\frac{1}{M}+3\right) = 16\varkappa^2\frac{\alpha+2\beta}{\alpha},
\]
and consequently by (13)
\[
\frac{16\varkappa^2}{M} = \left(\frac{1}{M}-1\right)^3\left(\frac{1}{M}+3\right),
\]
or, arranged in order,
\weq{14}{\frac{1}{M^4} - \frac{6}{M^2} + \frac{8(1-2\varkappa^2)}{M} - 3 = 0.}
If we therefore express our formulae in terms of $u, v$, the result of this consideration is the following:

By the substitution
\[
z = -\frac{\xi[(v^2+2vu^3)+u^6\xi]^2}{[v^2+(u^6+2vu^3)\xi]^2}
\]
the transformation is effected
\[
\frac{dz}{\sqrt{z(1-z)(1-v^8z)}} = \frac{v+2u^3}{v}\,\frac{d\xi}{\sqrt{\xi(1-\xi)(1-u^8\xi)}},
\]
provided that the modular equation (12) holds between $u, v$.

For a given $u$ the modular equation yields four values of $v$, and hence four different transformations of the third degree.
